% =====================================================================
\chapter{Gradient Descent}
% =====================================================================

\labinfo{03\_Gradient\_Descent}{gradient\_descent.py}{UCI Concrete Compressive Strength: 1,030 concrete
mixes, 8 inputs}{Regression (linear model trained by gradient descent)}{Test $R^2=0.835$, RMSE 6.66 MPa;
SGD lands on the exact least-squares weights}

\section{Dataset}
The \textbf{Concrete Compressive Strength} dataset (UCI, Yeh 1998) has 1,030 concrete mixes. Seven
inputs give the amount of each ingredient in kg per m$^3$ of mix (cement, blast-furnace slag, fly ash,
water, superplasticizer, coarse aggregate, fine aggregate), and the eighth is the age of the sample in
days (1 to 365). The target is the compressive strength in MPa, from 2.3 to 82.6 (mean 35.8).

Concrete gains most of its strength in the first few weeks and then levels off, so strength grows
roughly with the \emph{logarithm} of age. I therefore used $\log(\text{age})$ instead of age; this
single change is what allows a linear model to do well here. The inputs were standardised with the
training statistics, and the data was split 80/20 (824 training, 206 test mixes).

\section{Code and experiments}
\begin{itemize}
\item \code{mse\_loss} and \code{mse\_gradient} are the two formulas above, written with matrix products.
\item \code{gradient\_descent(X, y, lr, epochs, batch\_size)} is a single function for all three
      variants. \code{batch\_size=None} means batch GD, \code{32} means mini-batch, \code{1} means SGD.
      It shuffles the rows each epoch, applies the update for every batch, records the loss after
      each epoch, and stops early if the loss becomes infinite.
\item The script also computes the exact least-squares weights $\vect w^*$ with the normal equation and
      the largest eigenvalue of $X^\top X/n$, so the gradient-descent results can be checked against
      the theory.
\end{itemize}
\noindent The experiments are: (1)~the three variants for 200 epochs each; (2)~batch GD with five
learning rates, the last one just above the theoretical limit; (3)~a picture of the path that gradient
descent takes on a two-weight version of the problem.

\section{Results}

The largest eigenvalue of $X^\top X/n$ was $2.199$, so theory says batch GD is stable only for
$\eta<2/2.199=0.909$.

\begin{table}[H]
\centering\footnotesize\setlength{\tabcolsep}{4.5pt}
\caption{The three variants after 200 epochs. $\|\vect w-\vect w^*\|$ is the distance to the exact
least-squares weights. Times are from my run and will differ on other machines.}
\label{tab:gd-variants}
\begin{tabular}{@{}l r r S[table-format=2.3] S[table-format=1.3] S[table-format=1.3] S[table-format=1.4] S[table-format=1.3]@{}}
\toprule
\textbf{Method} & \textbf{Updates} & \textbf{Time} & {\textbf{Train $J$}} & {\textbf{Test MAE}} & {\textbf{Test RMSE}} & {\textbf{Test $R^2$}} & {$\|\vect w-\vect w^*\|$}\\
\midrule
Batch GD, $\eta=0.1$          & 200     & 0.01 s & 27.000 & 5.374 & 6.840 & 0.8258 & 7.449\\
Mini-batch (32), $\eta=0.01$  & 5,200   & 0.05 s & 26.246 & 5.242 & 6.697 & 0.8330 & 2.740\\
SGD, $\eta=0.001$             & 164,800 & 1.1 s  & 26.131 & 5.182 & 6.656 & 0.8350 & 0.100\\
\midrule
Normal equation (exact)       & --      & --     & 26.128 & 5.184 & 6.659 & 0.8349 & 0.000\\
\bottomrule
\end{tabular}
\end{table}

All three methods end up with almost the same test error as the exact solution (test $R^2$ between 0.826
and 0.835). SGD came within $0.1$ of the exact weights. Batch GD, after its 200 updates, still has a
loss of 27.0 instead of 26.13, but its weights are surprisingly far away (distance 7.4) even though its
predictions are nearly as good. The reason is that the seven ingredient amounts are tied to each other: they
describe one cubic metre of concrete, so their total is almost the same for every mix (about
2,340~kg, varying by only about 3\%). The loss surface therefore has a long, flat
valley in which very different weight vectors give almost the same predictions, and gradient descent
crawls along that valley slowly.

\begin{table}[H]
\centering\small
\caption{Batch gradient descent with different learning rates (200 epochs, exact minimum $J=26.128$).}
\label{tab:gd-lr}
\begin{tabular}{@{}l c c c c c@{}}
\toprule
$\eta$ & 0.001 & 0.01 & 0.1 & 0.5 & 0.955 \\
\midrule
Final training loss & 525.6 & 44.49 & 27.00 & 26.13 & diverged \\
\bottomrule
\end{tabular}
\end{table}

\begin{figure}[H]
\centering
\includegraphics[width=0.48\linewidth]{\figpath{03\_Gradient\_Descent}{gd_variants_loss.pdf}}\hfill
\includegraphics[width=0.48\linewidth]{\figpath{03\_Gradient\_Descent}{learning_rate_effect.pdf}}
\caption{Loss curves for GD variants and learning-rate stability.}
\end{figure}

\section{Source code}
\lstinputlisting[caption={\texttt{03\_Gradient\_Descent/gradient\_descent.py}}]{03_Gradient_Descent/gradient_descent.py}
\clearpage
